\documentclass[12pt]{article}
\usepackage[utf8]{inputenc}
\usepackage{amsmath, amssymb, graphicx, xcolor, nicefrac, hyperref}
\usepackage{float, multirow}
\usepackage{geometry}
\usepackage{textcomp}
\geometry{margin=1in}

\title{Empirically reported heritable fitness variance predicts severe depletion of neutral diversity in wild populations}
\author{Walid Mawass}
\date{}

\begin{document}
\section{Introduction}
Additive genetic variance in relative fitness, $V_A$, is the immediate currency of adaptive change: under Fisher’s fundamental theorem, the partial change in mean fitness attributable strictly to natural selection equals $V_A$ per generation on the mean-standardized scale, saying nothing on subsequent generations (Fisher 1930; Price 1972). Its magnitude bounds both the contemporary rate of adaptation and the capacity for evolutionary rescue (Crow 1958; Lande 1982; Gomulkiewicz and Holt 1995), making its estimation in natural populations a long-standing objective.

In an idealized neutral Wright–Fisher population, family sizes follow a Poisson distribution and the variance effective size approximately equals the census adult size, $N_e \approx N$ (Wright 1938). Because high-fitness parents transmit the alleles responsible for their selective advantage, their genetic contributions compound across generations. This heritable variance overdisperses long-term family sizes, accelerates the drift of neutral alleles, and depresses $N_e$ below the number of breeding adults, $N$ (Robertson 1961; Buffalo and Coop 2020). For a neutral site assorting independently of every selected locus in an otherwise Poisson background, Santiago and Caballero (1995, 1998) derived the closed form:

\begin{equation}
\frac{N_e}{N} = \frac{1}{1 + Q^2 V_A},
\end{equation}
where $Q = 2$ ($Q^2 = 4$) represents the cumulative transmission factor across generations, and $V_A$ corresponds to their mean-standardized additive variance, $C^2$. Expressed as the reciprocal, $N/N_e = 1 + 4 V_A$ scales linearly with $V_A$. Unlike linked background selection—which reduces diversity exponentially—the unlinked Robertson effect can impose only a modest reduction for realistic single-generation variances. Because pairwise nucleotide diversity reflects long-term coalescence ($\pi = 4 N_e \mu$), any sustained $V_A$ implies an upper bound on genome-wide neutral diversity, which in wild populations remains orders of magnitude below census expectations (Lewontin 1974; Leffler et al. 2012; Buffalo 2021). While this framework has been extended into a quantitative-genetic theory of linked background selection (Santiago and Caballero 2016; Buffalo and Kern 2024), the genome-wide unlinked Robertson reduction has largely been omitted from empirical accounting (Matheson and Masel 2025).

The pedigree-based animal model made $V_A$ estimable in unmanipulated wild populations by using deep monitoring records to separate additive genetic variance from environmental and maternal confounding (Henderson 1975; Kruuk 2004; Wilson et al. 2010). Applied to composite lifetime breeding success in long-term vertebrate systems—including St. Kilda Soay sheep and Wytham Woods great tits—it estimates additive variance in relative fitness directly (Postma 2014). Syntheses across these systems report substantial heritable variance across wild taxa, with estimates of order $V_A \sim 0.1$ (Bonnet et al. 2022).

Yet a sustained variance of that magnitude must also be mechanistically maintained. At mutation–selection–drift balance, standing additive variance in relative fitness is determined by the influx of new mutations and the distribution of fitness effects (DFE); when purifying selection dominates genetic drift ($2 N_e s_{\text{het}} \gg 1$), this expectation collapses to $V_A \approx U_d \bar{s}_{\text{het}}$ (Crow 1970; Barton 1990; Charlesworth 2015), where $U_d$ is the diploid genome-wide deleterious mutation rate and $\bar{s}_{\text{het}}$ is the mean heterozygous selection coefficient. Both parameters are now independently measured in vertebrates through pedigree sequencing (Bergeron et al. 2023) and DFE inference (Eyre-Walker and Keightley 2007; Huber et al. 2017). There is, consequently, an established mutational expectation for standing fitness variance that is derived entirely independently of neutral diversity.

Two distinct theoretical frameworks thus constrain the same biological quantity—one through neutral coalescent diversity, the other through mutation–selection balance. To our knowledge, the large fitness variances inferred from pedigree animal models have not been evaluated against either constraint. This is a one-to-one comparison: animal-model estimates represent observed-scale, mean-standardized additive variance in relative fitness ($I_A$)—the exact currency of Santiago and Caballero's $C^2$ and mutation–selection balance models alike. However, while both theoretical frameworks govern the long-run equilibrium variance, animal models estimate a local, single-generation realization conditioned on contemporary environmental states. Evaluating measured variances against these models tests what contemporary estimates would imply were they sustained across evolutionary time.

Here, we formulate an integrated analytical framework that couples the Robertson baseline—corrected for extra-pair paternity and mating systems—with linked selection integrated across empirical recombination maps and functional genomic annotations. Inverting this model yields the maximum additive fitness variance compatible with observed levels of neutral diversity. Applying this dual constraint to published pedigree estimates across monitored wild vertebrates reveals a systematic discordance among contemporary quantitative-genetic estimates, neutral diversity, and the mutational capacity of vertebrate genomes.

\section{Theory and Model}
\subsection{The quantity, and the two equilibria that constrain it}
The central quantity is the additive genetic variance in relative fitness, $V_A$: the variance of individual additive genetic breeding values for lifetime reproductive contribution, standardized to mean fitness ($\bar{w} = 1$). This corresponds to the mean-standardized additive variance $I_A$ (Crow 1958; Houle 1992) and matches Santiago and Caballero's (1995) $C^2$, providing a unified currency across all models evaluated below.

Both theoretical constraints require the long-run equilibrium value of this variance. Because pairwise nucleotide diversity ($\pi = 4 N_e \mu$) reflects neutral genealogies with an expected coalescent time of $2N_e$ generations (and a total pairwise branch length of $4N_e$ generations), standing diversity integrates the history of fitness variance over thousands of generations. Similarly, mutation–selection–drift balance defines the stationary distribution of standing variance under recurrent mutation and purging, rather than a transient, single-generation realization.

These dynamics operate across two distinct persistence windows:
\begin{enumerate}
    \item \textbf{The transmission memory ($Q$):} For a neutral site assorting independently of all selected loci, the inherited parental breeding value halves each generation due to Mendelian segregation. The cumulative transmission factor,
    \begin{equation}
    Q = \sum_{t=0}^{\infty} \left(\frac{1}{2}\right)^t = 2 \implies Q^2 = 4,
    \end{equation}
    is dominated by the earliest generational steps ($>85\%$ achieved within two generations). Consequently, establishing the unlinked Robertson baseline requires only that the direction of selective advantages remains consistent across roughly two consecutive generations—a condition satisfied under mutation–selection–drift balance, directional sweeps, and temporally fluctuating selection provided environmental fluctuations are positively autocorrelated across generations ($\tau_{\text{env}} \ge 2$).
    \item \textbf{The coalescent memory ($\pi$):} In contrast, the timescale required for neutral diversity to reach mutation–drift equilibrium following a shift in variance is categorically longer, requiring on the order of $2N_e$ to $4N_e$ generations.
\end{enumerate}

Connecting these frameworks requires equating two definitions of effective population size. Santiago and Caballero's derivation yields the asymptotic variance effective size, $N_{eV}$, which quantifies the per-generation rate of neutral drift resulting from the inherited compounding of family-size overdispersion. Nucleotide diversity measures the coalescent effective size, $N_{eC}$, which determines the timescale of the neutral genealogy. Under long-run demographic and selective equilibrium, the asymptotic rates of drift, inbreeding, and coalescence converge ($N_{eV} = N_{eC} = N_e$), providing the mathematical bridge that allows neutral diversity to constrain the standing additive variance in fitness. Key notation is summarized in Table~1.

\subsection{The Unlinked Robertson Baseline}
In an idealized diploid Wright--Fisher population of constant census size $N$, family size follows a binomial distribution that converges to Poisson in large populations ($V_k \approx \bar{k} = 2$), yielding an effective size approximately equal to the census adult size, $N_e \approx N$ (Wright 1938; Crow and Kimura 1970). This neutral benchmark assumes that variance in reproductive success is entirely non-heritable: each parent's genetic contribution represents an independent draw, generating no persistent fitness differences among genealogical lineages.

The Robertson effect (Robertson 1961) formalizes the consequence of violating that assumption. When reproductive success has an additive genetic basis, high-fitness parents transmit their favorable alleles to their offspring, establishing positive multi-generational autocovariance in lineage reproductive success (Buffalo and Coop 2020). In a population of constant size, the sustained compounding of high-fitness lineages accelerates the extinction of low-fitness lineages. This inflates the long-term variance of ancestral genetic contributions far beyond the Poisson expectation, depressing the asymptotic variance effective size, $N_e$, below the number of breeding adults.

For a neutral allele assorting independently of all selected loci ($r = 1/2$), Appendix~A integrates these selective-variance contributions across generations alongside baseline neutral sampling drift. Substituting the cumulative transmission factor established above ($Q^2 = 4$) yields the Santiago and Caballero (1995) unlinked baseline:
\begin{equation}
\frac{N_e}{N} = \frac{1}{1 + 4V_A}, \label{unlinked_baseline}
\end{equation}
where $V_A$ is the mean-standardized additive genetic variance in relative fitness ($I_A \equiv C^2$), evaluated in an otherwise Poisson demographic background.

Equation~\eqref{unlinked_baseline} carries an important scope condition. Because it depends strictly on the Mendelian dilution of transmitted breeding values under free assortment, it requires no assumptions regarding the genealogical persistence or haplotype structure of the selected background. It therefore applies across diverse selective regimes—including mutation-selection-drift balance, directional sweeps (with $V_A$ interpreted as the time-averaged variance, $\bar{V}_A$), and autocorrelated fluctuating selection—without requiring the stationarity parameters introduced below. Everything from \autoref{sec:linkage} onward requires explicit parameterization of haplotype longevity and recombination architecture; Equation~\eqref{unlinked_baseline} serves as an invariant, parameter-sparse ceiling on the effective population size.

\subsection{Variance Partitioning across Sexes}
To apply this framework to natural populations with separate sexes, we must account for the correlation of reproductive outcomes between mating partners. Maintaining a constant census size $N$ requires that each sex contributes, on average, two successful gametes per individual to the next generation ($\bar{k}_m = \bar{k}_f = 2$). The effective population size is determined by the pooled sampling variance of these gametic contributions (Hill 1979; Caballero 1994):

\begin{equation}
    N_e = \frac{4N}{V_{km} + V_{kf}}
\label{ne_sexes}
\end{equation}

Under strict monogamy, male and female reproductive outputs are perfectly correlated. The excess variance in family size generated by heritable fitness differences accrues jointly to both partners, maximally concentrating $V_A$ onto the combined gametic pool. To formalize this constraint, we partition the male gametic variance, $V_{km}$, into within-pair and extra-pair components as a function of the extra-pair paternity rate, $\alpha$. As $\alpha$ increases, extra-pair matings decouple male siring success from the female partner's reproductive output, allowing $V_{km}$ and $V_{kf}$ to accrue more independently.

Integrating this partitioned variance (Appendix B) modifies the unlinked baseline by a mating-system scalar, $\kappa(\alpha)$. The exact formulation of $\kappa(\alpha)$ depends on the covariance between extra-pair siring success and male breeding value. If extra-pair paternity is distributed independently of a male's additive genetic value for fitness, the scalar resolves to:

\begin{equation}
    \kappa(\alpha) = 1 + (1-\alpha)^2
\end{equation}

Alternatively, if extra-pair siring success is proportional to male additive genetic value, the scalar becomes:

\begin{equation}
    \kappa(\alpha) = 2(1 - \alpha + \alpha^2)
\end{equation}

Both formulations recover a maximal value of $\kappa = 2$ under strict monogamy, i.e. $\alpha = 0$, reflecting the two-fold amplification of heritable variance in fitness when reproductive outcomes are paired. Incorporating the mating-system scalar $\kappa(\alpha)$ into \autoref{unlinked_baseline} yields the unlinked expectation for the reduction in effective population size in a population with separate sexes:

\begin{equation}
    \frac{N_e}{N} = \frac{1}{1 + 4 \kappa(\alpha) V_A}
    \label{unlinked_mating}
\end{equation}

\subsection{Integration of Linked Constraint Across the Genetic Map}
The unlinked baseline assumes spatially uniform selective interference across the genome. Physical linkage breaks this symmetry by maintaining statistical associations between a neutral site and nearby selected loci across multiple generations, amplifying selective interference beyond the single-generation decay of independent assortment. We quantify this spatially structured amplification along a continuous genetic map using the framework of Santiago and Caballero (1995, 1998).

The cumulative selective interference on a neutral locus is governed by the temporal scaling factor $Q$: the sum of the per-generation genetic covariance between the neutral locus and the selected background, weighted by the joint survival probability of the physical association and the selective variance. The resulting reduction in effective population size is proportional to $Q^2 V_A$:

\begin{equation}
    \frac{N_e}{N} = \frac{1}{1 + Q^2 V_A}
    \label{linked_baseline}
\end{equation}

Here, $V_A$ represents the additive genetic variance in relative fitness scaled to unit mean ($\bar{w} = 1$; equivalent to $C^2$ in Equation 18 of Santiago and Caballero 1995).

Evaluating $Q^2$ across a continuous genetic map requires specifying $Z$, the per-generation survival rate of the selective variance. A defined $Z$ presupposes a stationary selected background, a condition satisfied under mutation-selection-drift equilibrium but violated by episodic selective sweeps, whose allele-frequency trajectories are inherently non-stationary (Barton 2000). By assuming $V_A$ reflects a steady-state balance between mutational input and purifying selection, this framework formalizes the reduction of neutral diversity as a continuous, quantitative-genetic manifestation of background selection (Santiago and Caballero 2016; Buffalo and Kern 2024).

Under this equilibrium regime, $Z = 1 - V_m/V_A$, where $V_m \approx U_d s^2$ is the per-generation mutational input to fitness variance, $U_d$ is the diploid deleterious mutation rate, and $s$ is the heterozygous selection coefficient. Under strong purifying selection, the House of Cards approximation—where allelic variation is governed by the balance between new mutations and selection—yields an equilibrium variance of $V_A \approx U_d s$ (Turelli 1984; Bürger 2000). The mutational input ratio therefore reduces to $V_m/V_A \approx s$, setting $Z \approx 1 - s$. In this regime, the lifetime of statistical associations is bounded by the selective timescale $\sim 1/s$ (Barton 1995; Santiago and Caballero 1998).

Conversely, under weak selection, individual mutational effects are small enough to maintain a Gaussian distribution of allelic effects, resulting in $V_A \gg V_m$. This yields $V_m/V_A \to 0$ and $Z \to 1$. Selective variance is effectively permanent, allowing physical associations between a neutral allele and its selected background to persist across many generations before dissipating. Biologically, this means the genetic background conferring high or low fitness remains intact long enough to force linked neutral alleles to rise or fall in tandem with it, driving severe local reductions in diversity. This sustained covariance forms the theoretical foundation for empirical methods that infer the genome-wide impact of background selection and selective sweeps by quantifying the localized depletion of neutral variation from genomic sequence data (e.g., McVicker et al. 2009; Elyashiv et al. 2016; Buffalo and Kern 2024).

Assuming heritable fitness variance is distributed uniformly across a chromosome of map length $M$ Morgans, the genetic distance $\ell = |x - y|$ between a selected locus $x$ and a neutral locus $y$ maps to the recombination fraction $r(\ell)$ via the Haldane function:

\begin{equation*}
    r(\ell) = \frac{1}{2}\left(1 - e^{-2\ell}\right)
\end{equation*}

We use the Haldane mapping function rather than the linear approximation ($r \approx \ell$). A linear formulation is adequate when calculating interference within short genomic segments, where the effect on $N_e$ is evaluated near the center of a narrow physical window (e.g., Buffalo and Kern 2024). However, because we integrate the interference of selected loci across the entirety of full chromosomes—many of which exceed 0.5 Morgans in our empirical datasets—a linear mapping would artificially inflate the contribution of distant sites. The Haldane function saturates at $r \to \frac{1}{2}$ as $\ell \to \infty$, ensuring that distantly linked sites on the same linkage group correctly approach independent assortment.

To calculate the joint probability that the association survives a single reproductive event, we multiply the probability that no crossover breaks the physical linkage, $(1 - r(\ell))$, by the probability that the selective variance persists, $Z$. Assuming these are independent events within a single generation, their product is $(1 - r(\ell))Z = \frac{Z}{2}(1 + e^{-2\ell})$. Integrating the infinite geometric sum of this joint probability over all selected loci (density $1/M$) yields the position-specific interference profile:
\begin{equation}
    Q^2(y) = \frac{1}{M} \int_0^M \left(\frac{1}{1 - \dfrac{Z}{2}\left(1 + e^{-2|x-y|}\right)}\right)^{2} dx
\end{equation}

Averaging $Q^2(y)$ over all neutral sites reduces the double integral to a single integral over pairwise genetic distances:
\begin{equation}
    \overline{Q^2}_{\text{linked}} = \frac{2}{M^2} \int_0^M (M - \ell) \left(\frac{1}{1 - \dfrac{Z}{2}\left(1 + e^{-2\ell}\right)}\right)^{2} d\ell
    \label{average_linked}
\end{equation}

This integral is evaluated numerically for each chromosome. For completeness, a closed-form approximation of this integral derived under the linear short-chromosome limit ($r \approx \ell$) is provided in Appendix A.

We approximate the reduction in effective population size across the entire chromosome by scaling the additive variance from the unlinked derivation by the amplification factor $\overline{Q^2}_{\text{linked}}$ and the mating-system scalar $\kappa(\alpha)$. We evaluate this reduction at the level of individual chromosomes because empirical linkage maps for these non-model organisms lack the fine-scale resolution required to model localized interference peaks. Because the local additive variance residing on a single chromosome, $v_i$, is a sufficiently small fraction of the genome-wide $V_A$, the discrete geometric decline is well approximated by its continuous exponential form:

\begin{equation}
    \left(\frac{N_e}{N}\right)_i \approx \exp\left( - \kappa(\alpha)\,\overline{Q^2}_{\text{linked}}\,v_i \right)
    \label{chrom_reduc}
\end{equation}

Propagating the strong ($Z = 1-s$) and weak ($Z \to 1$) selection regimes through the chromosomal limits of $\overline{Q^2}_{\text{linked}}$ establishes the theoretical boundaries of the linked interference (Table \ref{tab:limits}). In the limit of complete linkage ($M \to 0$), $\overline{Q^2}_{\text{linked}} \to (1-Z)^{-2}$. Under strong selection, this yields $\overline{Q^2}_{\text{linked}} \to s^{-2}$ and a local effective size of $\exp(-\kappa v_i / s^2)$. The reduction in neutral diversity is finite because the rapid purging of deleterious lineages limits their cumulative persistence. Under weak selection, $(1-Z)^{-2} \to \infty$; the near-permanent retention of selective variance locks neutral alleles to their genetic backgrounds across all generations. This drives local neutral diversity to zero, meaning all variation on the chromosome is eventually purged by background selection or fixed by genetic draft.

In the long-chromosome limit ($M \to \infty$), independent assortment dominates and $\overline{Q^2}_{\text{linked}} \to (1-\nicefrac{Z}{2})^{-2}$. Under strong selection, the amplification factor converges to $4/(1+s)^2$. Evaluating this for biologically realistic selection coefficients ($s \sim 0.01–0.10$) yields an amplification factor between 3.3 and 3.9. This demonstrates that even under strong selection, long chromosomes approach the maximal Robertson baseline of 4, bounding local diversity near $\exp(-4 \kappa v_i / (1+s)^2)$. Under weak selection, the amplification factor approaches 4, recovering the genome-wide unlinked baseline ($\nicefrac{N_e}{N} \to \exp(-4 \kappa v_i)$) regardless of the specific selection coefficient.

\begin{table}[H]
\centering
\caption{Limiting values of the amplification factor ($\overline{Q^2}_{\text{linked}}$) and the resulting local effective population size ($\nicefrac{N_e}{N}$) under strong ($Z = 1-s$) and weak ($Z \to 1$) purifying selection at the two extremes of chromosome map length ($M$). Here, $v_i$ represents the additive variance residing on the focal chromosome.}
\label{tab:limits}
\renewcommand{\arraystretch}{1.5}
\begin{tabular}{llcc}
\hline
\textbf{Regime} & \textbf{Metric} & \textbf{Complete Linkage} ($M \to 0$) & \textbf{Free Recombination} ($M \to \infty$) \\
\hline
\multirow{2}{*}{\textbf{Strong Selection} ($Z = 1-s$)}
& $\overline{Q^2}_{\text{linked}}$ & $s^{-2}$ & $\frac{4}{(1+s)^2}$ \\
& $\nicefrac{N_e}{N}$ & $\exp\left(-\frac{\kappa v_i}{s^2}\right)$ & $\exp\left(-\frac{4 \kappa v_i}{(1+s)^2}\right)$ \\[1ex]
\multirow{2}{*}{\textbf{Weak Selection} ($Z \to 1$)}
& $\overline{Q^2}_{\text{linked}}$ & $\infty$ & $4$ \\
& $\nicefrac{N_e}{N}$ & $0$ & $\exp\left(-4 \kappa v_i\right)$ \\
\hline
\end{tabular}

\end{table}

This structure establishes a critical asymmetry: the magnitude of linked interference is sensitive to the strength of purifying selection on short chromosomes, but converges toward the universal Robertson baseline on long chromosomes. For species with microchromosome-dominated karyotypes—such as avian taxa—a substantial fraction of the genome resides below the saturation scale of the Haldane mapping function. The predicted reduction in diversity on these microchromosomes is therefore highly sensitive to the parameterization of $U_d$ and $s$, whereas macrochromosomes are largely governed by the theoretical asymptote.

\subsection{Total Reduction in Diversity and Genomic Architecture}
The total, genome-wide reduction in effective population size combines selective interference from physically linked loci on the focal chromosome with the independent background variance segregating across the remainder of the genome. We partition the global additive variance $V_A$ among chromosomes in proportion to their functional target sizes. Assuming the majority of additive variance arises from functional sequence—an assumption we subsequently relax by incorporating conserved non-coding elements—the variance attributable to chromosome $i$ is $v_i = V_A f_i$, where $f_i$ is the fraction of the genome's total functional sequence. A focal neutral allele on chromosome $i$ is physically linked only to selected loci on its own haplotype, which carries variance $v_i/2$. Selected loci on the homologous haplotype and on all other chromosomes segregate independently of the focal allele. These loci contribute to the unlinked background, which carries the remaining variance $V_A - v_i/2$.

These two variance components enter the expected effective population size through distinct biological mechanisms. The linked component converges to an exponential function of the local linked load, $\exp(-\kappa (v_i/2)\,\overline{Q^2}_i)$, because it represents the continuous product of survival probabilities across the physical map of the chromosome. Conversely, the unlinked component accumulates through the aggregate interference of freely segregating loci across the rest of the genome, where independent assortment halves the statistical association between the focal neutral site and the selected background each generation, retaining the functional form of the Robertson geometric series (\autoref{unlinked_mating}). Assuming cross-chromosomal interference is negligible (a standard approximation because independent assortment prevents the sustained multi-generational buildup of linkage disequilibrium between distinct chromosomes), the two components act independently to yield the expected reduction in effective population size for a neutral allele on chromosome $i$:

\begin{equation}
    \frac{N_e}{N} \approx \frac{\exp\left(-\kappa\,\dfrac{v_i}{2}\,\overline{Q^2}_i\right)}{1 + 4\kappa\left(V_A - \dfrac{v_i}{2}\right)}
\label{total_red}
\end{equation}

The Relative Linkage Effect, $\Omega_i$, quantifies the proportional reduction in effective population size caused specifically by physical linkage. This isolates the effect of chromosomal architecture from the genome-wide baseline, mirroring recent advances that formally integrate linked and unlinked background selection in finite genomes (Matheson and Masel 2024). It is defined as the ratio of the total chromosome-specific reduction (\autoref{total_red}) to the global unlinked expectation (\autoref{unlinked_mating} evaluated at $V_A$):

\begin{align}
    \Omega_i &= \frac{(N_e/N)_{\text{full}}}{(N_e/N)_{\text{unlinked}}} \\[4pt]
    &= \exp\left(-\kappa\,\frac{v_i}{2}\,\overline{Q^2}_i\right)\left(\frac{1 + 4\kappa V_A}{1 + 4\kappa\left(V_A - \dfrac{v_i}{2}\right)}\right)
\label{linked_penalty}
\end{align}

Our theoretical framework directly predicts $N_e$, from which neutral diversity can be inferred. Because nucleotide diversity scales proportionally with effective size under the neutral coalescent ($\pi = 4N_e\mu$), $\Omega_i < 1$ implies that the physical architecture of chromosome $i$ depletes local neutral diversity beyond the unlinked expectation. We note that the second factor in \autoref{linked_penalty} is always greater than one: relocating $v_i/2$ from the global unlinked background into the linked term slightly reduces the genome-wide overdispersion of family sizes, partially relieving the unlinked Robertson interference. The net Linkage Effect, $\Omega_i$, reflects the balance between this minor global relief and the severe local interference imposed by physical linkage. This dynamic—where clustering selected sites on a chromosome intensifies local coalescent collapse while modestly reducing their global interference because they no longer segregate independently—is a recognized feature of multi-locus background selection models (Charlesworth et al. 1993).

To generate a genome-wide expectation of $N_e$ comparable to empirical measurements of nucleotide diversity, we weight the chromosome-specific predictions by each chromosome's physical length, $L_i$. This reflects the uniform probability that a random neutral mutation falls anywhere along the sequence. By restricting this integration to the autosomal sequence, we assume the contribution of sex chromosomes to the global variance pool is proportional to their physical length, an assumption we evaluate further in the Discussion.

\begin{equation}
    \mathbb{E}\left[\frac{N_e}{N}\right] = \sum_{i=1}^{n} \frac{L_i}{L_{\text{total}}} \cdot \frac{\exp\left(-\kappa \,\dfrac{v_i}{2}\,\overline{Q^2}_i\right)}{1 + 4\kappa\left(V_A - \dfrac{v_i}{2}\right)}
\end{equation}

The Global Linkage Effect, $\overline{\Omega}$, summarizes the aggregate impact of genomic architecture across the entire autosomal complement. Because the unlinked baseline is a genome-wide constant, $\overline{\Omega}$ reduces to the length-weighted mean of the individual chromosomal effects:
\begin{equation}
    \overline{\Omega} = \sum_{i=1}^{n} \frac{L_i}{L_{\text{total}}}\,\Omega_i
\end{equation}

A value of $\overline{\Omega} = 1$ indicates that the physical organization of the genome—chromosome lengths, recombination rates, and the spatial distribution of selected loci—imposes no interference beyond independent assortment. Values below one quantify the additional reduction in genome-wide effective size caused by physical linkage, providing a single standardized metric of how genomic architecture amplifies the total selective load across species.

\subsection{The Upper Bound of Heritable Variance in Relative Fitness}
The framework developed above predicts the genome-wide reduction in effective population size for a given empirical $V_A$. We invert this mapping to derive the maximum additive genetic variance in relative fitness compatible with maintaining a specified fraction of neutral diversity, a diagnostic limit defined as $V_{A,\text{max}}$.

Let $\delta \in (0,1)$ denote this minimum retained fraction of diversity. This threshold is equivalent to a target ratio of $N_e/N$, inferred under the standard coalescent assumption that the observed level of genome-wide nucleotide diversity scales proportionally with the effective population size. The unlinked baseline (\autoref{unlinked_mating}) provides the foundation for this theoretical limit. Because any degree of physical linkage amplifies selective interference ($\Omega_i \leq 1$), the unlinked expectation defines the absolute upper boundary of achievable diversity for any given $V_A$. No arrangement of recombination rates or functional architecture can restore diversity above this baseline. Consequently, $V_{A,\text{max}}$ derived from the unlinked model represents the most biologically permissive scenario; the reality of linked chromosomes dictates that a population will experience an even steeper reduction in effective population size at a given level of heritable variance. Setting the unlinked expectation equal to the threshold $\delta$ and solving directly for the maximum variance yields:

\begin{equation}
    V_{A,\text{max}} = \frac{1 - \delta}{4\kappa \delta}
\end{equation}

This limit depends exclusively on the specified diversity threshold $\delta$ and the mating-system scalar $\kappa$. Obligate monogamy ($\kappa = 2$) halves the sustainable $V_{A,\text{max}}$ relative to a randomly mating population with separate sexes ($\kappa = 1$). This reduction reflects the severe coalescent consequences of social pairing: by locking reproductive outcomes between partners, monogamy concentrates the heritable component of family-size variance simultaneously across both the male and female gametic pools, halving the magnitude of $V_A$ required to substantially deplete genome-wide diversity.

If a pedigree-based animal-model estimate of $V_A$ in a monitored wild population exceeds $V_{A,\text{max}}$ evaluated at the population's empirically observed $N_e/N$, the ecological measurements are irreconcilable with the expectations of a rigorous theoretical model at stable mutation-selection-drift equilibrium. In such cases, at least one of the following biological or methodological conditions must hold: (i) the population is not at equilibrium and is undergoing an active, transient loss of neutral variation; (ii) the true effective breeding population substantially exceeds the monitored census via cryptic immigration from unsampled, unstructured demes; or (iii) the statistical fit of the animal model captures non-additive or environmentally structured covariance among relatives.

Because the animal model detects statistical associations between phenotypic similarity and relatedness without identifying the biological mechanism, it is highly susceptible to absorbing indirect genetic effects, localized kin-structured interactions, and spatially autocorrelated environments into the $V_A$ component (Kruuk and Hadfield 2007; Stopher et al. 2012). Furthermore, variance component estimation is sensitive to pedigree structure (Mawass and Milot 2025). These factors inflate the statistical estimate of $V_A$ in the pedigree without generating the proportional, transmissible reproductive advantage required to compress the coalescent timescale, and by extension, deplete observable neutral diversity—two measures that theoretically and empirically track each other closely in genome-wide data (Charlesworth 2009; Buffalo 2021).

\subsection{The Architecture of Variance: Infinitesimal and Oligogenic Limits}
The continuous integration derived above distributes the chromosome-specific selective variance, $v_i = V_A f_i$ (scaled by its functional target size), uniformly across the genetic map of chromosome $i$. This assumption aligns with the classical infinitesimal model of quantitative genetics, where $V_A$ arises from an infinite number of loci, each of negligible effect. However, the true genetic architecture of fitness in natural populations remains empirically unresolved (Barton and Keightley 2002). To bound the potential outcomes, we must evaluate the continuum from fully infinitesimal to oligogenic, where a small number of loci contribute substantial variance. Because the linked interference profile $Q^2(y)$ varies systematically with the genetic position $y$ of the focal neutral site, the aggregate reduction in diversity depends strongly on the spatial distribution of these causal loci.

To formalize this dependence, let $g(x) \geq 0$ define an arbitrary density of selective variance along chromosome $i$, normalized such that $\int_0^M g(x)\,dx =1$. The selective variance per unit map length at the selected locus position $x$ is $v_i g(x)/2$. The position-specific interference profile at a neutral site $y$ becomes:
\begin{equation}
    Q^2(y,\,g) = \int_0^M g(x)\left(\frac{1}{1-\dfrac{Z}{2}\left(1 + e^{-2|x-y|}\right)}\right)^2dx
\end{equation}

Averaging over all neutral sites under the uniform density $1/M$ and exchanging the order of integration via Fubini's theorem yields:

\begin{align}
    \overline{Q^2}(g) &= \int_0^M g(x)\left[\frac{1}{M}\int_0^M K(x,y)\,dy\right]dx \\
    &= \int_0^M g(x)\,Q^2(x)\,dx \label{general_g}
\end{align}
here, the symmetry of the kernel $K(x,y)$ allows the inner integral to resolve to $Q^2(x)$. \autoref{general_g} establishes that the chromosome-averaged interference is simply the $g$-weighted average of the spatial interference profile $Q^2(x)$.

Under the infinitesimal model ($g(x) = \nicefrac{1}{M}$), \autoref{general_g} recovers $\overline{Q^2}_i$ as derived previously (\autoref{average_linked}). Under an oligogenic model comprising $J$ loci of equal effect at positions $x_1,\dots, x_J$, the density is $g(x)= J^{-1} \sum_{j=1}^J \delta (x - x_j)$, and the chromosome-averaged interference resolves to a discrete mean:

\begin{equation}
    \overline{Q^2}(g) = \frac{1}{J}\sum_{j=1}^J Q^2(x_j) \label{oligogenic_mean}
\end{equation}

The magnitude of the linked reduction therefore depends directly on whether causal loci occupy regions of high or low $Q^2(x)$. Because $Q^2(x)$ is maximized at the chromosome midpoint $M/2$ and minimized at the telomeres—a consequence of centrally located loci being physically linked to sequence on both chromosome arms, whereas telomeric loci are linked to only one—centrally located causal loci contribute disproportionately to the chromosome-averaged reduction in diversity.

The implications for model bias are governed by the convexity of the exponential linked reduction term, $\exp(-\kappa v_i \,\overline{Q^2}(g)/2)$. Because the second derivative of this function with respect to $\overline{Q^2}(g)$ is positive, Jensen's inequality bounds the expectation for a single oligogenic locus at a uniformly random position $x_0 \sim \text{Uniform}[0,M]$:

\begin{equation}
    \mathbb{E}_{x_0}\left[\exp\left(-\frac{\kappa v_i}{2}\,Q^2(x_0)\right)\right] > \exp\left(-\frac{\kappa v_i}{2}\,\mathbb{E}[Q^2(x_0)]\right) = \exp\left(-\frac{\kappa v_i}{2}\,\overline{Q^2}_i\right)
\end{equation}

Biologically, because the interference effect scales exponentially, placing a selected locus in a region of high interference (the chromosome center) reduces $N_e$ far more than placing it in a region of low interference (the telomere) relieves it. Consequently, the average effective population size produced by a randomly placed causal locus is always greater than the effective size produced by spreading that same variance uniformly across the chromosome. This inequality constitutes a formal proof that the continuous infinitesimal model is conservative when evaluated against randomly positioned oligogenic variation. The magnitude of this Jensen's gap is proportional to $(\kappa v_i/2)^2 \cdot \text{Var}_x[Q^2(x)]$. Thus, the infinitesimal model most underestimates $N_e$ on chromosomes carrying substantial selective variance and exhibiting steep spatial gradients in $Q^2(x)$.

However, this conservatism assumes causal loci are positioned randomly with respect to the recombination landscape. If causal loci are systematically enriched in regions of reduced recombination—such as pericentromeric domains or inversions, where empirical quantitative trait loci are frequently mapped (Kirkpatrick and Barton 2006; Yeaman 2013)—the local interference $Q^2(x_j)$ systematically exceeds the chromosome-wide average $\overline{Q^2}_i$. Such non-random spatial enrichment rapidly closes the Jensen's gap and can reverse it, rendering the infinitesimal model an underestimate of the actual reduction in effective population size.

While the Jensen's inequality framework evaluates unobservable causal loci, the integration in \autoref{general_g} accommodates empirical functional densities in place of theoretical uniform distributions. At the chromosome level, using annotated protein-coding gene counts, $G_i$, to define the functional partition coefficients ($\hat{f}_i = G_i / G_{\text{total}}$) provides only a minor refinement over using physical coding sequence length, as the two metrics are highly correlated across vertebrate karyotypes. The critical architectural uncertainty lies within the chromosome. For species with integrated physical-to-genetic maps, the spatial distribution of selective variance can be localized to the specific genetic positions of annotated genes, $x_j^{(i)}$. Substituting this empirical architecture into \autoref{oligogenic_mean} yields:

\begin{equation}
    \overline{Q^2}_i\left(\hat{g}_i\right) = \frac{1}{G_i}\sum_{j=1}^{G_i} Q^2\left(x_j^{(i)}\right)
\end{equation}

This substitution replaces the assumption of uniform variance per unit genetic distance with uniform variance per coding gene. While this parameterization resolves the spatial clustering of functional targets along the recombination map, it does not identify which specific genes harbor the variance estimated by the animal model. The unobservable identity of the causal loci remains the fundamental source of theoretical uncertainty, governing the ultimate direction and magnitude of the Jensen's gap.

\section{Materials and Methods}
\subsection{Empirical Pedigree Parameters and Variance Posteriors}
Estimates of additive genetic variance in relative fitness ($V_A$) were compiled from published pedigree-based animal model analyses of longitudinally monitored wild vertebrate populations. Studies were included if they: (i) reported $V_A$ for a composite fitness measure or a major fitness component (annual or lifetime reproductive success, survival to reproductive age, or a multivariate index thereof); (ii) provided sufficient pedigree depth and sample size to support reliable separation of additive from residual environmental variance; and (iii) reported posterior credible intervals or standard errors for the $V_A$ estimate. All estimates were expressed as additive genetic variance in fitness on the relative fitness scale ($\bar{w} = 1$), consistent with the parameterization of Equations (1) and (5); published estimates reported on a different scale were standardized prior to analysis [author: describe standardization procedure]. Where multiple published estimates were available for the same population, [author: state resolution rule—most recent, most comprehensive pedigree, or population-level average]. The full dataset, with source studies and population identifiers, is provided in [Table/Supplementary Table X].

\subsection{Mating System Parameterization}
For each species, the extra-pair paternity rate $\alpha$ was obtained from published molecular parentage analyses of the same or a conspecific monitored population [author: specify source for each species and confirm whether $\alpha$ is defined as the proportion of offspring or the proportion of broods containing at least one extra-pair offspring; the theoretical definition in Equations (3) and (4) is the proportion of offspring, and any conversion should be noted here]. For species lacking a direct estimate of $\alpha$, [author: state whether strict monogamy was assumed or whether those species were excluded from the mating-system-corrected prediction].
The mating-system scalar $\kappa(\alpha)$ was computed under both the random extra-pair siring model (Equation 3) and the fitness-proportional extra-pair siring model (Equation 4), and predictions are reported under both to bracket the effect of the assumed extra-pair mechanism.

\subsection{Genomic Architecture and Functional Target Size}
For species with chromosome-level genome assemblies, reference sequences and gene annotation files (GTF or GFF format) were retrieved from NCBI. Analysis was restricted to fully assembled autosomes with unambiguous correspondence to a major linkage group in the published linkage map (see below); unplaced scaffolds, unlocalized contigs, sex chromosomes, and mitochondrial sequence were excluded. Each assembled autosome was assigned to its corresponding linkage group on the basis of marker content and ordering reported in the linkage map source publication.

The between-chromosome partition of additive fitness variance was computed from annotated protein-coding gene counts following Equation (19):
\begin{equation*}
    \hat{f}i = \frac{G_i}{G_{\text{total}}}
\end{equation*}
where $G_i$ is the number of annotated protein-coding loci on autosome $i$ and $G_{\text{total}} = \sum_i G_i$ over all included autosomes. Gene counts were taken from the feature tables of NCBI assembly annotation reports. For species without available genome assemblies, $f_i$ was approximated as proportional to physical chromosome length, $f_i \approx L_i / L_{\text{total}}$.

For the mutational influx calculation (below), total genome-wide coding sequence length $L_{\text{coding}}$ was estimated independently of the gene-count partition. All CDS features in the annotation file were extracted and overlapping intervals on each strand collapsed into non-redundant contiguous ranges to prevent multiple isoforms and overlapping reading frames from inflating the functional target; the genome-wide collapsed coding length is $L_{\text{coding}} = \sum_i L_{\text{coding},i}$. This quantity enters only the $U_d$ calculation and does not affect the $\hat{f}_i$ partition.

\subsection{Genetic Map Lengths}
Genetic map lengths ($M_i$, in Morgans) were taken from published empirical linkage maps for each species. Where sex-specific and sex-averaged map lengths were both available, the sex-averaged value was used to reflect the expected map length across both parental gametes. Autosomes without a resolved major linkage group in the published map were excluded from the linked-constraint calculation; the selective variance apportioned to excluded chromosomes by $\hat{f}_i$ was redistributed to the unlinked background [author: verify this is the procedure applied, or describe the alternative].

For species lacking per-chromosome map data, $M_i$ was approximated as $M_i = \bar{\rho} \times L_i \times 10^{-8}$, where $\bar{\rho}$ is the published genome-wide average recombination rate in cM per Mb and $L_i$ is the physical chromosome length in base pairs.

The closed-form expressions for $Q^2(y)$ and $\overline{Q^2}_i$ (Equations 8 and 9) were derived under the linear approximation $r \approx |x - y|$, which is accurate for $|x - y| \ll 0.5$ Morgan. For the vertebrate chromosomes analyzed here, the dominant contribution to $\overline{Q^2}_i$ arises from closely linked locus pairs where this approximation is most accurate. For chromosomes substantially longer than 0.5 Morgan, the linear approximation overestimates the recombination fraction between distant loci relative to the Haldane mapping function, causing a mild underestimate of $Q^2$ from those pairs; because this produces a conservative (upward) bias in the predicted $N_e/N$, the approximation errs in the direction of understating constraint. The magnitude of this bias across the chromosome-length distribution of each species is quantified in the Supplementary Material.

\subsection{Mutational Parameters and the Mutation-Selection Balance Test}
The per-generation survival rate of selective variance, $Z = 1 - V_m / V_A$, requires an estimate of the mutational input $V_m \approx U_d s^2$ (Equations [X] and [X] in Theory). The genome-wide diploid deleterious mutation rate was estimated as:
\begin{equation*}
    U_d = 4\mu L_{\text{coding}}
\end{equation*}
where $\mu$ is the published per-base-pair per-generation point mutation rate for each species and the factor of four converts from a per-haplotype per-base-pair rate to a diploid rate over both coding and regulatory sequence. The regulatory target is approximated as equal in length to the annotated exome, following empirical estimates of the ratio of constrained non-coding to coding sequence in vertebrate genomes [cite: e.g., Meader et al. 2010 or equivalent]; this multiplier doubles the effective functional target relative to coding sequence alone and enters the calculation as the factor of two in $U_d = 2\mu \times 2L_{\text{coding}}$.
Because the distribution of fitness effects of deleterious mutations is not independently constrained for the non-model vertebrate populations studied here, we do not fix a point estimate of $s$. Instead, the implied average heterozygous selection coefficient $s_d$ was inferred by solving $V_m = U_d s_d^2$ for $s_d$ given the empirically estimated $V_m = V_A(1 - Z)$ at each value of $V_A$. This was evaluated across the joint 95\% credible intervals of the published $V_A$ posterior and the empirical uncertainty in $\mu$, enabling a test of whether the maintenance of observed standing variance requires selection coefficients consistent with known vertebrate DFE parameters and tolerable levels of mutational load.

\subsection{Empirical Nucleotide Diversity}
\textit{[Author input required. This subsection must describe: (i) the source of genome-wide or synonymous-site nucleotide diversity $\pi$ for each species in the dataset, including the specific studies, variant call sets, or databases consulted; (ii) the genomic compartment used for the diversity estimate (e.g., fourfold-degenerate sites, synonymous sites, intergenic sequence, or genome-wide); (iii) the operational definition of $N_e/N$ used for empirical comparison—whether computed as $\pi/(4\mu)$ divided by census size $N$, or derived from a directly reported $N_e$ estimate; and (iv) how census size $N$ was estimated or obtained for each monitored population, since $N_e/N$ is the predicted quantity and $N$ enters the denominator. Without this subsection the theoretical predictions in Equations (11) and (13) cannot be connected to any empirical observation, which is the central comparison of the paper.]}

\section{Results}



\subsection{Intragenomic Variation in Reduction}
While the expected species-wide diversity, $\mathbb{E}[N_e/N]$, provides a global summary of evolutionary constraint, the continuous spatial framework predicts highly heterogeneous topographies of diversity within the genome. By evaluating the Main Equation across independent chromosomes, we can analytically partition the sources of intragenomic variance.

For a focal chromosome $i$, the genome-wide unlinked baseline generated by independent assortment is captured by the denominator of the Main Equation: $1 + 4\kappa(V_A - v_i/2)$. Because the variance segregating on a single haplotype represents a minor fraction of the genome-wide total or $v_i/2 \ll V_A$, this unlinked constraint is effectively uniform across the genome. Consequently, any predicted variance in standing neutral diversity between chromosomes is driven strictly by the localized, constraint of physical linkage.

Because the localized selective variance is proportional to the functional target size of the chromosome, i.e. $v_i = V_A f_i$, the survival of neutral lineages is governed by the structural ratio of functional density to recombinational map length. Chromosomes with high coding densities and restricted recombination suffer severe exponential decay in local $N_e$, driving rapid lineage sorting. Conversely, chromosomes with expansive map lengths and sparse functional targets effectively isolate neutral alleles from selective backgrounds. These structural dynamics predict the existence of macro-genomic regions that maintain robust neutral diversity even when the global, genome-wide constraint is strong.


\section{Discussion}
The synthesis of quantitative genetics and coalescent theory places a fundamental limit on the heritable fitness variance segregating in wild populations: ecological estimates from contemporary pedigrees must be reconcilable with the neutral diversity maintained over deep evolutionary timescales. Integrating the Robertson effect with empirical mating-system structure, continuous chromosomal architecture, and annotation-derived functional target sizes demonstrates that current empirical estimates of $V_A$ routinely violate this limit. Across the [N\_SPECIES] wild vertebrate species in our dataset, the unlinked baseline yields predicted $N_e/N$ ratios from [XX] to [XX] (median [XX]) across the posterior distributions of reported $V_A$. Incorporating the mating-system scalar $\kappa(\alpha)$ and chromosome-specific selective interference further depresses these expectations to genome-wide $N_e/N$ values between [XX] and [XX]. For [X\_OUT\_OF\_N] species, the reported $V_A$ exceeds $V_{A,\text{max}}$, the theoretical ceiling of variance compatible with observed sequence variation. Conversely, based on the standard coalescent expectation that neutral diversity scales proportionally with effective population size ($\pi = 4N_e\mu$), these species maintain empirical genome-wide nucleotide diversity between [XX] and [XX]. This translates to sequence-derived $N_e/N$ estimates of [XX]–[XX] (Corbett-Detig et al. 2015; Buffalo 2021; Lewin and Eyre-Walker 2026)—values two to three orders of magnitude above our theoretical predictions. The magnitude of this mismatch between ecologically estimated fitness variance and coalescent-scale diversity forms a quantitative paradox that cannot be resolved by parametric uncertainty alone.

The magnitude of this discrepancy persists even under permissive parameterizations of the genomic and mutational inputs. Expanding the functional target size to include a conserved regulatory component equivalent in length to the annotated exome, and evaluating the diploid deleterious mutation rate $U_d$ across the full range of empirical uncertainty, maximizes the potential for mutational turnover to replenish fitness variance. Despite this, the per-generation variance survival rate ($Z = 1 - V_m/V_A$) approaches its theoretical maximum across all species. Because the mutational influx to fitness variance ($V_m \approx U_d s^2$) is bounded by empirical $U_d$ and realistic selection coefficients, the massive $V_A$ estimated from contemporary pedigrees necessarily drives $Z \to 1$. Biologically, $Z \to 1$ indicates that statistical associations between neutral alleles and their selected backgrounds persist across many generations before dissipating via recombination or independent assortment. This sustained covariance amplifies the chromosome-wide interference ($\overline{Q^2}_\text{linked}$) toward its ceiling, concentrating the coalescent histories of neutral alleles into the reproductive trajectories of the linked selected sites.

Furthermore, because the Relative Linkage Effect ($\Omega_i$) is bounded at or below one, no spatial arrangement of selected loci can restore neutral diversity above the global unlinked expectation. Physical linkage can only exacerbate the interference imposed by independent assortment; it cannot relieve it. The Global Linkage Effect ($\overline{\Omega}$) ranges from [XX] to [XX] across species, confirming that chromosomal architecture universally amplifies the unlinked reduction in $N_e$, with the most severe interference localized to gene-dense chromosomes possessing short genetic map lengths.

This dynamic is starkly evident in the avian taxa: the bimodal avian karyotype sequesters a substantial fraction of functional targets on microchromosomes. Because these microchromosomes frequently exhibit map lengths well below 0.5 Morgans—the scale at which the Haldane mapping function fails to saturate and interference approaches the limit of complete linkage ($\overline{Q^2}_\text{linked} \to (1-Z)^{-2}$)—the predicted reduction in local diversity becomes highly sensitive to the parameterization of $Z$, and by extension, the assumed distribution of fitness effects for deleterious mutations. The inability of bounding parameterizations to resolve this mismatch in any species establishes a fundamental incompatibility between the magnitude of ecologically estimated fitness variance and the mechanics of the neutral coalescent.

This incompatibility bears directly on Lewontin’s paradox: the observation that genome-wide nucleotide diversity varies across taxa by far less than their census population sizes would predict, implying a systematic, non-linear decoupling of effective and census sizes (Lewontin 1974). Extensive comparative genomic surveys confirm that $N_e$ inferred from nucleotide diversity varies by roughly three orders of magnitude, whereas census sizes span more than eight (Leffler et al. 2012; Buffalo 2021). Linked selection—comprising background selection against recurrent deleterious mutations and interference from selective sweeps—is the consensus proximate mechanism suppressing $N_e$ relative to $N$. Furthermore, the magnitude of this reduction ($N_e/N$) correlates positively with species-level proxies for the total intensity of linked selection (Corbett-Detig et al. 2015).

The theoretical framework presented here formalizes this dynamic by specifying the coalescent mechanism—heritable variance in fitness amplifying the Robertson effect across both unlinked and physically linked architectures—that drives the reduction in $N_e/N$. However, substituting pedigree-derived variance estimates into this framework reverses the standard theoretical gap. While classical background selection models frequently underestimate the reduction in $N_e$ required to match observed diversity (Corbett-Detig et al. 2015; Comeron 2017), our model parameterized with empirical $V_A$ predicts reductions two to three orders of magnitude more severe than those observed. Although the predicted coalescent mechanism correctly suppresses diversity, the magnitude of fitness variance estimated by animal models in the wild is mathematically irreconcilable with the deep evolutionary maintenance of nucleotide variation. This structural mismatch restates Lewontin’s paradox in quantitative genetic terms, establishing the transmissible component of fitness variance as a deeply constrained observable that must jointly satisfy both contemporary phenotypic models and historical coalescent diversity.

The structural mismatch between contemporary variance estimates and historical diversity strongly suggests that a substantial fraction of the $V_A$ estimated by pedigree-based animal models in wild populations does not constitute stable, transmissible additive genetic variance. Linear mixed models structured by pedigree relatedness detect statistical associations between phenotypic similarity and genealogical proximity. However, high statistical fit does not imply biological validity in partitioning transmission genetics (Kruuk 2004; Hadfield et al. 2010). Because related individuals in structured wild populations frequently experience spatially or socially autocorrelated environments—such as high-quality territories, maternal provisioning, or localized resource patches linked to natal dispersal—common-environment effects generate phenotypic covariance among relatives that the basic animal model erroneously absorbs into the additive genetic component (Kruuk and Hadfield 2007; Stopher et al. 2012). In kin-structured populations, indirect genetic effects—where an individual's phenotype depends on the genotypes of social partners—further confound this separation, producing substantial upward bias in $V_A$ estimates when unmodeled (Baud et al. 2022; Mawass and Milot 2025).

These biases are most acute for fitness traits, as realized reproductive success is phenotypically correlated with nearly all life-history traits sensitive to environmental quality. Furthermore, models using unweighted lifetime reproductive success (LRS) as a fitness proxy systematically overestimate Fisherian additive variance; age-structured demographic models establish that LRS fails to appropriately discount delayed reproduction and generation-time variation, artificially inflating variance estimates for iteroparous vertebrates (Sæther and Engen 2015; de Villemereuil et al. 2020). The shallow, unbalanced pedigrees typical of wild populations also lack the structural power required to separate additive from dominance and epistatic contributions, allowing non-additive variance to inflate $V_A$ estimates (Wolak and Keller 2014; [Insert citation to your JEB paper here]). This instability is empirically supported by temporal fluctuations in $V_A$ estimates across years within the same monitored populations, a pattern incompatible with a stable additive architecture [CITE temporal instability of $V_A$ results for specific monitored populations in dataset].

The coalescent framework developed here quantifies this bias. By mapping the theoretical ceiling of sustainable variance ($V_{A,\text{max}}$) directly to observed nucleotide diversity, it provides a strict upper bound on the stable, transmissible variance a population can harbor without undergoing coalescent collapse. The fraction of reported $V_A$ that exceeds $V_{A,\text{max}}$ ranges from [XX]\% to [XX]\% across the analyzed species. This excess defines the minimum non-transmissible fraction of the reported variance, indicating that the statistical signal driving these animal models is fundamentally disconnected from the evolutionary mechanics governing long-term molecular diversity.

Alternatively, this discrepancy may arise from violations of the demographic assumptions inherent in the coalescent framework rather than statistical bias in the ecological variance estimates. The monitored populations in our dataset are defined by the practical geography of fieldwork—nest box grids, island boundaries, or marked cohorts—rather than the biological boundaries of the effective breeding unit. Because wild study populations are typically demes within larger, structured metapopulations, cryptic immigration from unsampled conspecifics continuously replenishes local variation. If this immigration occurs at rates comparable to or exceeding the local coalescent rate, it maintains diversity even when local selective interference suppresses the within-deme $N_e$. Theoretical work on the interaction of gene flow and linked selection establishes that the restoration of diversity depends on the ratio of immigration to local drift and selection, with modest gene flow capable of maintaining diversity well above the local selective equilibrium (Whitlock and Barton 1997; Charlesworth et al. 2003).

Gene flow at rate $m$ per generation introduces neutral haplotypes with distinct selective histories, driving local diversity toward the metapopulation average over a timescale of approximately $1/m$ generations. For the passerine populations most heavily represented in the quantitative genetics literature, natal dispersal rates suggest $m$ on the order of [XX] per generation. This implies a restoration timescale of [XX] generations—long relative to the pedigree observation window, but short relative to the coalescent timescale governing the maintenance of $\pi$. The capacity of this demographic replenishment to maintain observed $\pi$ against the predicted local reduction in $N_e/N$ depends directly on metapopulation geometry, dispersal distance distributions, and the spatial scale of the environmental covariance underlying $V_A$. Future extensions of this framework must couple spatially explicit metapopulation models with the Robertson effect to determine the threshold at which locally estimated $V_A$ ceases to predict species-wide sequence variation.

A final resolution involves violations of the stationarity assumption underlying the derivation of the continuous interference profile $\overline{Q^2}_\text{linked}$. The chromosome-wide constraint assumes that $V_A$ remains approximately constant across generations—a condition satisfied under mutation-selection-drift equilibrium on a stable fitness landscape. This assumption fails if the direction or magnitude of selection fluctuates. Temporally variable selection attenuates the trans-generational autocovariance in relative fitness that drives the Robertson effect: lineages expressing high fitness in one generation may be selected against in the next as the optimum shifts. This oscillation disrupts the multi-generational association between neutral alleles and their selected genetic backgrounds, collapsing the $Q^2$ amplification of selective interference (Buffalo and Coop 2019). The critical requirement for this attenuation is that fitness reversals occur on a timescale comparable to or shorter than $1/s$ generations, where $s$ is the mean fitness differential; reversals operating on longer timescales permit substantial coalescent compression before selection alters direction.

Rapidly oscillating environmental optima, interannual variation in favored genotypes, frequency-dependent dynamics, and shifting spatial fitness gradients can all generate this required attenuation in wild populations. Genomic evidence tracking temporal allele-frequency changes confirms the operation of selective interference at genome-wide scales, yet the observed magnitudes fall well below the expectations generated by pedigree-derived $V_A$ (Buffalo and Coop 2020). If temporally fluctuating selection drives this paradox, it demonstrates that pedigree-based $V_A$ estimated over short observation windows systematically overstates the temporal mean selective autocovariance governing the neutral coalescent. The relevant parameter for coalescent predictions is not the instantaneous statistical $V_A$ isolated in a single generation, but the long-run, time-averaged transmissible component. Because these three resolutions—pedigree estimation bias, metapopulation gene flow, and temporally fluctuating selection—are not mutually exclusive, identifying their relative contributions remains a major empirical challenge. By defining the theoretical limits of sustainable variance, the $V_{A,\text{max}}$ diagnostic framework provides a quantitative baseline to structure this empirical program.

The implied selection coefficients (Appendix D) demonstrate that contemporary $V_A$ estimates are incompatible with steady-state mutation-selection-drift balance. Under the House of Cards approximation—the strong-selection limit where new mutations exert large effects relative to pre-existing standing variation—the equilibrium variance $V_A \approx U_d s$ defines the required average heterozygous selection coefficient as $s^* = V_A/U_d$. Given reported $V_A \in [0.10, 0.30]$ and empirical diploid deleterious mutation rates of $U_d \approx [XX]$ for the analyzed vertebrates, this relationship necessitates $s^* \approx [XX]–[XX]$. This expectation exceeds empirical estimates of the mean heterozygous selection coefficient by one to two orders of magnitude, as vertebrate distributions of fitness effects consistently cluster around $\bar{s} \approx 10^{-3}–10^{-2}$ (Eyre-Walker and Keightley 2007; Huber et al. 2017). The Robertson formula compounds this implausibility: it treats $V_A$ as a sufficient statistic for the unlinked interference effect, but this equivalence holds only in the small-$t$ limit where heterozygous selection coefficients are well below one. When fitness variance is concentrated in alleles of large heterozygous effect—precisely the regime implied by $\bar{t}^* = V_A/U_d$ at the reported $V_A$ magnitudes—the per-locus interference at unlinked sites scales as 2$t/(1+t)^2$ rather than $t$, so that strongly deleterious alleles contribute less genome-wide interference per unit variance than the Robertson formula predicts (Charlesworth 2012; Rettelbach et al. 2025). As shown in Appendix D, this DFE correction partially relieves the Robertson-predicted suppression at large $\bar{t}^*$ but does not resolve the paradox: the corrected $N_e/N$ predictions remain orders of magnitude below observation for all species in the dataset.

The weak-selection regime is similarly incompatible. Solving for the implied selection coefficient under the near-neutral approximation ($V_A \approx 2N_e U_d s^2$) and substituting the resultant values yields $N_e s \gg 1$ for all populations in our dataset (Appendix D). Because $N_e s \gg 1$ indicates that directional selection overwhelms genetic drift, this outcome invalidates the foundational premise of the weak-selection limit. Consequently, neither strong-selection nor weak-selection boundaries can recover the reported $V_A$ magnitudes when constrained by empirically calibrated mutation rates and selection coefficients. This inability to reconcile pedigree-derived variance with steady-state mutational limits establishes an independent, demographic-free analogue to the coalescent paradox.

The theoretical framework developed here relies on several structural approximations that define the boundaries of its application. First, the equilibrium derivation (Appendix D) assumes a uniform heterozygous selection coefficient ($s$) across all deleterious loci. Empirical distributions of fitness effects (DFE) in vertebrates, however, are strongly leptokurtic, characterized by a bulk of mildly deleterious mutations and a long tail of large-effect alleles. Because the per-locus contribution to additive variance scales proportionally with $s$ under House of Cards balance, large-effect alleles dominate $V_A$ (Eyre-Walker and Keightley 2007). Consequently, under an empirical DFE, the effective selection coefficient driving the Robertson effect exceeds the DFE mean. This dynamic renders the required selection coefficients ($s^*$) implied by pedigree $V_A$ even more biologically improbable than the uniform-$s$ approximation suggests, indicating that this architectural assumption underestimates the severity of the paradox. Second, the analytical framework models discrete, non-overlapping generations. The wild vertebrate populations analyzed here are iteroparous, necessitating age-structured demographic models that modify the mapping between family-size variance and $N_e$ (Engen et al. 2005). While this age-structure correction is moderate for species with low interannual survival variance, it must be integrated into future demographic modeling of specific iteroparous study systems. Third, the current formulation is restricted to autosomal inheritance. For avian taxa harboring substantial functional content on the Z chromosome, the hemizygous Z in females is transmitted without recombination, placing it in the complete-linkage limit (map length $M = 0$ in the heterogametic sex). Excluding sex chromosomes therefore systematically underestimates the total genome-wide selective interference for these species. Fourth, the spatial distribution of fitness variance within chromosomes is modeled as uniform per Morgan, conforming to the infinitesimal model. As established mathematically via Jensen’s inequality, distributing variance uniformly minimizes local interference compared to randomly positioned large-effect loci, generating an upper bound on $N_e/N$. However, if unobserved causal loci cluster in regions of suppressed recombination, this approximation reverses, underestimating the local interference. Because the empirical genomic locations of loci contributing to pedigree-derived $V_A$ remain unmapped for these monitored populations, this architectural uncertainty imposes a fundamental limit on the precision of spatial interference predictions. Finally, the predictions assume demographic and environmental stationarity. Populations experiencing recent bottlenecks, expansions, or habitat fragmentation violate this assumption. Because several monitored populations in our dataset occupy restricted island habitats or fragmented ranges, their current molecular diversity ($\pi$) reflects non-equilibrium historical demography. In such systems, resolving the fraction of the $N_e$ reduction attributable to historical bottlenecks versus contemporary selective interference requires independent demographic inference prior to applying the equilibrium framework.

This study establishes a quantitative bridge between the heritable fitness variance estimated from contemporary pedigrees and the neutral diversity predicted by coalescent theory, demonstrating that these two bodies of empirical data are mutually irreconcilable. By integrating the Robertson effect with mating-system structure, continuous chromosomal architecture, and empirical functional target sizes, this framework derives the theoretical maximum additive genetic variance ($V_{A,\text{max}}$) compatible with maintaining observed nucleotide diversity. The observation that reported $V_A$ exceeds this threshold by factors of [XX]–[XX] across [X\_OUT\_OF\_N] species constitutes a structural incompatibility that persists across all bounding parameterizations. Consequently, pedigree-based animal models likely capture a conflated mixture of stable Mendelian transmission and environmentally structured covariance among relatives. This non-genetic covariance inflates statistical estimates of $V_A$ but fails to generate the multi-generational selective interference necessary to reduce long-term $N_e$. Metapopulation gene flow and temporally fluctuating selection further decouple instantaneous phenotypic variance from long-term coalescent dynamics. Ultimately, these findings establish a fundamental limit on the interpretation of quantitative genetic parameters in wild populations: heritable variance in relative fitness cannot be treated as an isolated ecological metric decoupled from its population-genomic context. The magnitude of variance derived from pedigrees must be evaluated against the consequent reduction in effective population size before it can be interpreted as the stable, transmissible Mendelian variance that drives long-term evolutionary adaptation.

\appendix
\section{Appendix A: Derivation of the Unlinked Baseline}
To derive the baseline reduction in effective population size generated by unlinked heritable fitness variance, we track the cumulative variance in allele frequency change driven by the selected genetic background, independent of continuous spatial linkage.

In a neutral Wright-Fisher population with a constant diploid census size $N$, the variance in allele frequency change per generation due to random genetic drift for a neutral allele at frequency $p$ is the standard gametic sampling variance:

\begin{equation}
    \text{Var}(\Delta p^{\text{drift}}) = \frac{p(1-p)}{2N}
\end{equation}

When variance in reproductive success is heritable ($V_A > 0$), the neutral allele experiences selective interference from the genetic background upon which it resides. In a given generation, a diploid population of $N$ individuals has relative fitnesses $w_i$ with expectation $\mathbb{E}[w_i] = 1$ and variance $\text{Var}(w_i) = V_W$. The frequency of the neutral allele in the subsequent generation is approximately:

\begin{equation}
    p' \approx p + \frac{1}{2N}\sum_{i=1}^{N}(w_i - 1)z_i
\end{equation}
where $z_i \in \{0, \frac{1}{2}, 1\}$ denotes the allele copy number in individual $i$ under Hardy-Weinberg proportions. Because the neutral locus is unlinked from all selected sites, relative fitness $w_i$ and the neutral genotype $z_i$ are uncorrelated in expectation ($\mathbb{E}[(w_i-1)z_i] = 0$). The variance in allele frequency change driven purely by this single generation of heritable fitness variation is:

\begin{equation}
    \text{Var}(\Delta p^{\text{sel}}) =\frac{1}{4N^2}\sum_{i=1}^{N}\text{Var}((w_i - 1)z_i) \approx V_A \frac{p(1-p)}{2N}
\end{equation}

This formulation (Robertson 1961) relies on a large-$N$ approximation, assumes Hardy-Weinberg equilibrium ($\text{Var}(z_i) = p(1-p)/2$), and isolates the additive genetic component $V_A$ from the total fitness variance, as contributions from dominance and epistasis do not transmit stably across generations to drive multi-generational allele frequency changes.

Because the neutral allele and the selected background are unlinked (segregating on non-homologous chromosomes), they assort independently during meiosis ($r = 0.5$). Consequently, the statistical association, or linkage disequilibrium, between the neutral allele and its initial selected background halves each generation. The directional change in the neutral allele frequency driven by this specific selective cohort decays geometrically:
\begin{equation}
    \Delta p_t^{\text{sel}} = \left(\frac{1}{2}\right)^t \Delta p^{\text{sel}}
\end{equation}

Integrating this effect over time yields the total cumulative change in the neutral allele frequency resulting from a single generation of heritable variance:

\begin{equation}
    \Delta p_{\text{total}}^{\text{sel}} = \sum_{t=0}^{\infty} \left(\frac{1}{2}\right)^t \Delta p^{\text{sel}} = 2 \Delta p^{\text{sel}}
\end{equation}

Assuming $V_A$ remains stationary under mutation-selection-drift equilibrium, a new influx of selective variance of magnitude $V_A p(1-p)/(2N)$ is generated each generation. The long-run footprint of each generation's selective variation resolves to $2\Delta p^{\text{sel}}$.

The total selective variance in allele frequency change is the variance of this cumulative sum. Because the selective association persists across multiple generations, the variance of the cumulative effect scales with the square of the cumulative sum:
\begin{equation}
    \text{Var}(\Delta p_{\text{total}}^{\text{sel}}) = \text{Var}(2 \Delta p^{\text{sel}}) = 4 V_A \frac{p(1-p)}{2N}
\end{equation}
This establishes that the multi-generational persistence of associations across independently assorting chromosomes amplifies the effective variance four-fold.

Because the neutral allele is physically unlinked from the causal selected loci, the variance generated by selective interference and the variance generated by neutral demographic sampling arise from independent biological processes. Their contributions to the total variance in allele frequency are therefore additive. Equating the standard definition of effective gametic sampling variance to the sum of these independent variances yields:
\begin{equation}
    \frac{p(1-p)}{2N_e} = \text{Var}(\Delta p^{\text{drift}}) + \text{Var}(\Delta p_{\text{total}}^{\text{sel}})
\end{equation}

Substituting the expanded terms:
\begin{equation}
    \frac{p(1-p)}{2N_e} = \frac{p(1-p)}{2N} + 4 V_A \frac{p(1-p)}{2N}
\end{equation}

Dividing by the common neutral variance term $p(1-p)$ and multiplying by 2 yields the unlinked expectation for the effective population size:
\begin{equation}
    \frac{1}{N_e} = \frac{1}{N} (1 + 4V_A) \Rightarrow N_e = \frac{N}{1 + 4V_A}
\end{equation}
This baseline establishes the minimum selective interference generated by heritable fitness variance, irrespective of spatial chromosomal architecture. The extension to separate sexes and structured mating systems is derived in Appendix B, where the additive variance $V_A$ is scaled by the mating system parameter $\kappa(\alpha)$, yielding $N_e/N = 1/(1+4\kappa V_A)$.

\section{Appendix B: Variance Partitioning and the Mating System Scalar ($\kappa$)}

To quantify how mating system structure modifies the Robertson effect, we must partition the heritable variance in reproductive success between male and female gametic contributions. Because alleles pass alternately through male and female life histories, sex-specific differences in the mapping of additive genetic quality to realized offspring number dictate the total selective interference exerted on the neutral genome.

We evaluate a diploid, outcrossing population with a constant census size $N$, an equal sex ratio, and discrete generations. To maintain stationary demography, the mean family size for both males and females is $\bar{k}_m = \bar{k}_f = 2$.

Let female reproductive success be defined as $k_f = 2 + x_f + \epsilon_f$. Here, $x_f$ represents the heritable breeding value for fitness with variance $\text{Var}(x_f) = V_A$, and $\epsilon_f$ captures stochastic environmental noise yielding a neutral Poisson demographic baseline. Consequently, the female gametic contribution to selective interference is governed entirely by $V_A$. We assume that total pair fecundity is biologically limited by female quality (e.g., through maternal provisioning or egg production). The male benefits from his social mate's quality via within-pair paternity, but does not contribute an independent additive genetic effect to the total clutch size. This assumption holds broadly for socially monogamous passerines and is relaxed below in Model B to allow male genetic quality to drive extra-pair siring success.

Male reproductive success, $k_m$, comprises within-pair (WP) offspring sired with his social mate and extra-pair (EP) offspring sired with other females in the population. Let $\alpha$ represent the population-wide rate of extra-pair paternity. The male's within-pair success is directly coupled to his social mate's breeding value, $x_{f,\,\text{pair}}$, proportionally scaled by the fraction of the clutch he successfully defends, $1-\alpha$.

\subsubsection{Model A: Random Extra-Pair Paternity}
In Model A, extra-pair siring success operates as a demographic lottery, entirely decoupled from the male's own additive genetic quality. This condition applies when extra-pair copulations are driven by random spatial encounters rather than female choice for "good genes."

Under this dynamic, the selective component of a male's reproductive success relies exclusively on the fraction of his social mate's reproductive output that he secures:
\begin{equation}
    x_{m,\,\text{selective}} = (1-\alpha)x_{f,\,\text{pair}}
\end{equation}

Because scaling a random variable by a constant scales its variance by the square of that constant, the selective variance in the male gametic contribution is:
\begin{equation}
    \text{Var}(x_{m,\,\text{selective}}) = (1-\alpha)^2 \text{Var}(x_f) = (1-\alpha)^2 V_A
\end{equation}

The total effective selective variance driving allele frequency change in the population is the sum of the independent female and male components. Factoring out $V_A$ defines the mating system multiplier $\kappa$:
\begin{equation}
    V_{\text{total, selective}} = V_A + (1-\alpha)^2 V_A = \left[ 1 + (1-\alpha)^2 \right] V_A
\end{equation}

Thus, under random extra-pair mating:
\begin{equation}
    \kappa(\alpha) = 1 + (1-\alpha)^2
\end{equation}

\subsubsection{Model B: Selective Extra-Pair Paternity}
In Model B, females selectively allocate extra-pair copulations to males with high heritable fitness. We formalize this by defining a male's extra-pair siring success as directly proportional to his own breeding value, $x_m$, bounded by the population-wide availability of EP offspring, $\alpha$.

The selective component of a male's reproductive success is therefore the sum of his defended within-pair success and his earned extra-pair success:
\begin{equation}
    x_{m,\,\text{selective}} = (1-\alpha)x_{f,\,\text{pair}} + \alpha x_m
\end{equation}

To calculate the variance of this sum, we must specify the covariance between the male's breeding value, $x_m$, and his social mate's breeding value, $x_{f,\,\text{pair}}$. Assuming random social pair formation (no assortative mating for fitness), this covariance is strictly zero. The variance therefore expands additively:
\begin{equation}
    \text{Var}(x_{m,\,\text{selective}}) = (1-\alpha)^2 \text{Var}(x_f) + \alpha^2 \text{Var}(x_m)
\end{equation}

Assuming the underlying additive genetic variance for fitness is symmetric across sexes ($\text{Var}(x_f) = \text{Var}(x_m) = V_A$), the male selective variance resolves to:

\begin{equation}
    \text{Var}(x_{m,\,\text{selective}}) = \left( 1 - 2\alpha + \alpha^2 + \alpha^2 \right) V_A = \left( 1 - 2\alpha + 2\alpha^2 \right) V_A
\end{equation}

Summing the female and male selective gametic contributions yields the multiplier $\kappa$ for Model B:
\begin{equation}
    V_{\text{total, selective}} = V_A + \left( 1 - 2\alpha + 2\alpha^2 \right) V_A = \left( 2 - 2\alpha + 2\alpha^2 \right) V_A
\end{equation}

Factoring out the constant establishes the final scalar:
\begin{equation}
    \kappa(\alpha) = 2(1 - \alpha + \alpha^2)
\end{equation}

Evaluating the limits of Model B confirms its biological boundaries. Under strict social monogamy ($\alpha = 0$), $\kappa = 2$, reflecting that the female's quality perfectly dictates both her own and her mate's success. At the opposite boundary of complete extra-pair siring ($\alpha = 1$, functionally equivalent to a lekking system where social pairing does not constrain paternity), the model mathematically recovers $\kappa = 2$. Biologically, full extra-pair paternity shifts the entire male selective variance directly onto male genetic quality, perfectly restoring the total selective variance to the monogamous baseline.

The scalar reaches its minimum at $\alpha = 0.5$ ($\kappa = 1.5$). At this precise intermediate frequency, male reproductive success is partitioned equally between two independent genetic metrics (social mate quality and own quality). This orthogonal splitting maximizes the dilution of variance, minimizing the aggregate selective interference exerted on the coalescent.

The quantity $V_{\text{total, selective}} = \kappa(\alpha) V_A$ encapsulates the effective heritable variance in family size, unified across both sexes, that drives the Robertson geometric decay. Substituting $\kappa(\alpha) V_A$ into the unlinked baseline derivation yields the mating-system-corrected expectation for the reduction in effective population size:

\begin{equation}
    \frac{N_e}{N} = \frac{1}{1 + 4\kappa(\alpha) V_A}
\end{equation}
which corresponds to Equation (5) in the main text.

\section{Appendix C: Analytical Integration of the Chromosome-Wide Constraint}

This appendix derives the chromosome-wide average selective interference, $\overline{Q^2}_\text{linked}$, resulting from heritable fitness variance distributed continuously across a genetic map. The derivation proceeds in two stages. First, we formulate the integral under the exact Haldane mapping function—the expression we numerically evaluate to generate the predictions for every chromosome in our dataset. Second, we derive a closed-form analytical expression under the linear approximation ($r \approx \ell$), formally establishing its biological domain of validity and proving that it provides a strict lower bound on the magnitude of selective interference.

\subsection*{C.1\quad Position-specific constraint profile}

Consider a neutral locus located at genetic position $y$ on a chromosome of total map length $M$ Morgans. The cumulative selective interference exerted on this neutral site by a single locus under selection at position $x$ depends on the persistence of their statistical association. The joint probability that this association survives both recombination and the loss of selective variance across one generation is $(1 - r(\ell))Z$, where $\ell = \vert{}x-y\vert{}$ is the pairwise genetic distance.

The cumulative lifetime interference generated by this selected locus is the infinite geometric sum of this survival probability:

\begin{equation}
    \sum_{t=0}^{\infty}\left[(1 - r(\ell))\,Z\right]^t = \frac{1}{1-(1-r(\ell))Z}
\label{eq:geom_sum}
\end{equation}

Assuming selected loci are distributed uniformly across the genetic map (consistent with the infinitesimal model), we obtain the position-specific interference profile $Q^2(y)$ by squaring this cumulative effect and averaging over all possible positions of the selected locus, $x \in [0,M]$:

\begin{equation}
    Q^2(y) = \frac{1}{M}\int_0^M \left(\frac{1}{1-(1-r(|x-y|))Z}\right)^2 dx
\end{equation}

Under the Haldane mapping function, $r(\ell) = \frac{1}{2}(1-e^{-2\ell})$, the per-generation survival probability becomes:

\begin{equation}
    (1-r(\ell))\,Z = \frac{Z}{2}(1+e^{-2\ell})
\end{equation}

To resolve the absolute value in the genetic distance $\ell = \vert{}x-y\vert{}$, we split the spatial integral at the focal position $y$, computing the interference contributed by the chromosome arms to the left and right of the neutral site separately:

\begin{equation}
    Q^2(y) = \frac{1}{M}\left[\int_0^y f_\text{H}(\ell)\,d\ell + \int_0^{M-y} f_\text{H}(\ell)\,d\ell\right]
\end{equation}
where $f_\text{H}(\ell) = \left(1 - \frac{Z}{2}(1+e^{-2\ell})\right)^{-2}$. Because $f_\text{H}(\ell)$ is a transcendental function of the genetic distance $\ell$, the antiderivative cannot be expressed using elementary functions. Consequently, Equation (C.4) is evaluated numerically.

\subsection*{C.2\quad Chromosome-wide average under the Haldane mapping function}

To determine the aggregate interference experienced by neutral variation across the entire chromosome, we average the position-specific profile $Q^2(y)$ over all possible neutral-site positions, assuming a uniform physical density of neutral sites ($1/M$):
\begin{equation}
    \overline{Q^2}_\text{linked} = \frac{1}{M}\int_0^M Q^2(y)\,dy = \frac{1}{M^2}\int_0^M\int_0^M f_\text{H}(|x-y|)\,dx\,dy
\label{eq:double_int}
\end{equation}

We can reduce this double integral over the spatial domain $[0,M]^2$ to a single integral over all pairwise genetic distances $\ell = \vert{}x-y\vert{}$. Because the interference function is symmetric with respect to direction ($f_\text{H}(\vert{}x-y\vert{}) = f_\text{H}(\vert{}y-x\vert{})$), the contributions from loci where $x > y$ equal those where $x < y$:

\begin{equation}
\overline{Q^2}_\text{linked} = \frac{2}{M^2}\int_0^M\int_0^x f_\text{H}(x-y)\,dy\,dx
\end{equation}

Substituting $\ell = x - y$ in the inner integral ($dy = -d\ell$; moving the limits $y \in [0,x]$ to $\ell \in [x,0]$, which reverses the sign):

\begin{equation}
    = \frac{2}{M^2}\int_0^M\int_0^x f_\text{H}(\ell)\,d\ell\,dx
\end{equation}

The integration region is $\{(\ell,x) : 0 \leq \ell \leq x \leq M\}$. Exchanging the order of integration to evaluate the spatial limits before the genetic distance:

\begin{equation}
    \overline{Q^2}_\text{linked} = \frac{2}{M^2}\int_0^M f_\text{H}(\ell)\left(\int_\ell^M dx\right)d\ell = \frac{2}{M^2}\int_0^M (M-\ell)\,f_\text{H}(\ell)\,d\ell
\label{eq:haldane_avg}
\end{equation}

The biological intuition behind the linear weight $(M-\ell)$ is straightforward: on a chromosome of finite length $M$, the number of locus pairs separated by a given genetic distance $\ell$ decreases linearly as $\ell$ increases. As the genetic distance approaches the total length of the chromosome ($\ell \to M$), the number of physically possible pairs vanishes. This corresponds to Equation (10) in the main text.

\medskip
\noindent\textit{Verification of Biological Boundaries.} The integrity of the model requires that it recovers known biological limits. Using the identity $\frac{2}{M^2}\int_0^M(M-\ell)\,d\ell = 1$, we evaluate the boundaries:

\medskip
\noindent\textit{Complete-linkage limit} ($M \to 0$): For a chromosome with zero recombination (e.g., the avian W or mammalian Y chromosome), all pairwise genetic distances vanish ($\ell \to 0$). The survival function evaluates to $f_\text{H}(0) = (1-Z)^{-2}$ uniformly across the chromosome:

\begin{equation}
    \lim_{M\to 0}\overline{Q^2}_\text{linked} = (1-Z)^{-2}
\end{equation}

\noindent\textit{Independent-assortment limit} ($M \to \infty$): For an infinitely long chromosome, the Haldane function saturates at $r \to 1/2$. Consequently, $f_\text{H}(\ell) \to (1-Z/2)^{-2}$ for the vast majority of locus pairs. Because the exponential convergence to this asymptote is rapid, the dominated convergence theorem ensures:

\begin{equation}
    \lim_{M\to\infty}\overline{Q^2}_\text{linked} = \frac{1}{(1-Z/2)^2}
\end{equation}

These boundaries establish that as a chromosome becomes infinitely long, the selective interference precisely decays to the Robertson unlinked baseline derived in Appendix A ($Q^2 = 4$ as $Z \to 1$).

\subsection*{C.3\quad Closed-form approximation under the linear recombination model}

For genetic distances sufficiently short that interference is severe ($\ell \ll 0.5$ Morgans), we can approximate the Haldane function with a linear expansion $r(\ell) \approx \ell$. Under this linear mapping, the per-generation survival probability becomes $(1-\ell)Z$, yielding the integrand:

\begin{equation}
    f_\text{L}(\ell) = \left(\frac{1}{1-Z+\ell Z}\right)^2
\end{equation}

This approximation allows an analytical antiderivative. Using the substitution $u = 1-Z+\ell Z$, $du = Z\,d\ell$:

\begin{equation}
    \int_0^D f_\text{L}(\ell)\,d\ell = \frac{1}{Z}\int_{1-Z}^{1-Z+DZ}\frac{du}{u^2} = \frac{D}{(1-Z)(1-Z+DZ)}
\label{eq:antideriv}
\end{equation}

Applying this to the left and right chromosome arms ($D = y$ and $D = M-y$) yields the closed-form profile for a specific neutral site:

\begin{equation}
    Q^2_\text{L}(y) = \frac{1}{M(1-Z)}\left(\frac{y}{1-Z+yZ} + \frac{M-y}{1-Z+(M-y)Z}\right)
\end{equation}

To determine the chromosome-wide average, we integrate over all neutral-site positions:

\begin{equation}
    \overline{Q^2}_\text{L} = \frac{1}{M^2(1-Z)}\int_0^M\left(\frac{y}{1-Z+yZ} + \frac{M-y}{1-Z+(M-y)Z}\right)dy
\end{equation}

By symmetry, both arms contribute identically across the interval $[0,M]$, simplifying the integral to:
\begin{equation}
    \overline{Q^2}_\text{L} = \frac{2}{M^2(1-Z)}\int_0^M \frac{y}{1-Z+yZ}\,dy
\end{equation}

Using the substitution $u = 1-Z+yZ$, $dy = du/Z$, we evaluate the integral:

\begin{align}
    \int_0^M \frac{y}{1-Z+yZ}\,dy
    &= \frac{1}{Z^2}\int_{1-Z}^{1-Z+MZ}\left(1 - \frac{1-Z}{u}\right)du \\
    &= \frac{1}{Z^2}\left[u-(1-Z)\ln u\right]_{1-Z}^{1-Z+MZ}
\end{align}

Evaluating at the bounds and combining logarithmic terms yields:

\begin{equation}
    = \frac{M}{Z} - \frac{1-Z}{Z^2}\ln\left(1+\frac{MZ}{1-Z}\right)
\end{equation}

Substituting this back into the leading coefficient $2/(M^2(1-Z))$ provides the final closed-form expression for chromosome-wide interference under linear recombination:

\begin{equation}
    \overline{Q^2}_\text{L} = \frac{2}{MZ(1-Z)} - \frac{2}{M^2Z^2}\ln\left(1+\frac{MZ}{1-Z}\right)
\label{eq:closed_form}
\end{equation}

\subsection*{C.4\quad Domain of validity and directional bias of the approximation}

\noindent\textit{Agreement at complete linkage ($M \to 0$).} Setting $\epsilon = MZ/(1-Z)$ and applying a Taylor expansion to the logarithm ($\ln(1+\epsilon) = \epsilon - \epsilon^2/2 + O(\epsilon^3)$) yields:

\begin{align}
    \overline{Q^2}_\text{L} &= \frac{2}{MZ(1-Z)} - \frac{2}{M^2Z^2}\left[\frac{MZ}{1-Z} - \frac{1}{2}\left(\frac{MZ}{1-Z}\right)^2 + O(\epsilon^3)\right] \\
    &= \frac{1}{(1-Z)^2} + O(M)
\end{align}

This confirms that the linear approximation recovers the Haldane expectation in the limit of zero recombination.

\medskip
\noindent\textit{Failure at independent assortment ($M \to \infty$).} For highly recombining chromosomes, as $M \to \infty$, the leading term $2/(MZ(1-Z)) \to 0$ and the logarithmic correction decays even faster, driving $\overline{Q^2}_\text{L} \to 0$. This violates the biological reality established in (C.10): selective interference cannot decay below the unlinked baseline imposed by independent assortment. The linear approximation assumes recombination rates can exceed $0.5$, artificially purging associations that should persist across unlinked chromosomes.

\medskip
\noindent\textit{Proof of conservative bias.} Because the exponential decay of physical linkage slows as distance increases, the true recombination fraction is strictly concave ($e^{-2\ell} > 1-2\ell$ for all $\ell > 0$). Therefore, the exact Haldane mapping is always bounded above by the linear approximation:

\begin{equation}
    r_\text{H}(\ell) = \frac{1}{2}(1-e^{-2\ell}) < \ell = r_\text{L}(\ell)
\end{equation}

Because the linear approximation overestimates the true recombination rate for any non-zero distance, it systematically underestimates the per-generation survival probability of the statistical association between loci. Consequently, it underestimates the interference function pointwise across the entire chromosome ($f_\text{L}(\ell) < f_\text{H}(\ell)$).

Integrating this strictly lower function guarantees that:

\begin{equation}
    \overline{Q^2}_\text{L} < \overline{Q^2}_\text{H} \quad \text{for all } M > 0
\end{equation}

This establishes that the closed-form expression is a lower bound on the true severity of selective interference. When mapping this to the biological prediction, underestimating interference produces a conservative overestimate of the effective population size, or $N_e/N$.

\medskip
\noindent\textit{Application to empirical data.} The linear approximation is robust when the characteristic distance over which interference decays ($d^* \approx 1-Z$) is small compared to 0.5 Morgans. Biologically, this expression is highly reliable for avian microchromosomes ($M \lesssim 0.3–0.5$ Morgans), where nearly all pairwise locus interactions fall securely within the linear regime. However, for macrochromosomes spanning multiple Morgans, the discrepancy becomes substantial. Therefore, all primary predictions in the main text are derived via numerical integration of the exact Haldane formulation (Equation C.8), utilizing this closed-form expression strictly as an analytical cross-check for localized, short-range behavior.

\section{Appendix D: Stationary Equilibrium and Implied Selection Coefficients}
To determine whether the magnitude of additive genetic variance, $V_A$, estimated from contemporary pedigrees can be stably maintained under neutral sequence diversity, we analyze the balance of evolutionary forces at stationarity. At equilibrium, the per-generation influx of variance from novel mutation is exactly offset by the depletion of variance through purifying selection and genetic drift:

\begin{equation}
    \Delta V_{\text{mut}} + \Delta V_{\text{sel}} + \Delta V_{\text{drift}} = 0
\end{equation}

We derive each flux term, assuming a common additive heterozygous effect, $s$, across all loci, before extending the analysis to full empirical distributions of fitness effects (DFE).

\paragraph{Mutational influx.} In a population of size $N$, $U_d N$ new deleterious mutations occur each generation across the genome, where $U_d$ is the diploid deleterious mutation rate. A single new mutation arises at initial frequency $q_0 = 1/(2N)$. For an allele with additive heterozygous effect $s$, its contribution to additive variance at low frequency ($q \ll 1$) is $2q(1-q)s^2 \approx 2qs^2$. The initial variance contributed by a single new mutation is therefore $2q_0 s^2 = s^2/N$. Summing this contribution across all $U_d N$ new mutations yields the total mutational influx:
\begin{equation}
    \Delta V_{\text{mut}} = (U_d N) \cdot \frac{s^2}{N} = U_d s^2 \equiv V_m
\end{equation}

\paragraph{Depletion by purifying selection.} For a rare segregating deleterious allele ($q \ll 1$), the per-locus additive variance is $v \approx 2qs^2$. The per-generation change in frequency driven by additive directional selection is $\Delta q = -sq(1-q) \approx -sq$. The corresponding change in the per-locus variance is:

\begin{equation}
    \Delta v_{\text{sel}} = 2s^2 \Delta q = 2s^2(-sq) = -s(2qs^2) = -sv
\end{equation}

Summing the per-locus depletion across all segregating sites yields the total selective depletion of genome-wide variance:
\begin{equation}
    \Delta V_{\text{sel}} = -s V_A
\end{equation}

\paragraph{Depletion by genetic drift.} In a population of effective size $N_e$, stochastic gametic sampling erodes heterozygosity at a rate of $1/(2N_e)$ per generation ($\Delta H = -H/(2N_e)$). Because per-locus additive variance scales directly with heterozygosity ($v = s^2 H$), drift depletes $V_A$ at exactly the same rate:
\begin{equation}
    \Delta V_{\text{drift}} = -\frac{V_A}{2N_e}
\end{equation}

\paragraph{Equilibrium solution.} Setting the net flux to zero and solving for the stationary variance $V_A$:
\begin{equation}
    U_d s^2 - s V_A - \frac{V_A}{2N_e} = 0
    \quad \Rightarrow \quad
    V_A = \frac{2N_e U_d s^2}{1 + 2N_e s}
\label{eq:VA_MSD}
\end{equation}

This deterministic balance holds under the assumptions of rare alleles, additive effects, and constant population size. We verify its equivalence to classical continuous diffusion theory in Section D.4 below.

\subsection{The Implied Selection Coefficient}

By rearranging Equation (\ref{eq:VA_MSD}) as a quadratic in $s$:

\begin{equation}
    2N_e U_d s^2 - 2N_e V_A s - V_A = 0
\end{equation}

we isolate the positive root, determining the precise selection coefficient $s^*$ required to maintain a given empirical $V_A$ at stationary equilibrium:

\begin{equation}
    s^* = \frac{V_A + \sqrt{V_A^2 + 2U_d V_A/N_e}}{2U_d}
\label{eq:s_implied}
\end{equation}

This general expression bridges the strong-selection and near-neutral boundaries, allowing us to evaluate the biological plausibility of pedigree-derived variance estimates across all selective regimes.

\subsection{The Strong-Selection Boundary ($N_e s \gg 1$)}

When directional selection overwhelms genetic drift, the drift term $V_A/(2N_e)$ becomes negligible compared to selective depletion ($sV_A$). In this limit, Equation (\ref{eq:VA_MSD}) collapses to the classic House of Cards approximation (Turelli 1984; Bürger 2000):

\begin{equation}
    V_A \approx U_d s \qquad (N_e s \gg 1)
\end{equation}

Here, deterministic mutation-selection balance maintains deleterious alleles at frequency $q \approx \mu/s$. Solving for the implied selection coefficient:

\begin{equation}
    s^*_{\text{strong}} = \frac{V_A}{U_d}
\label{eq:sHoC}
\end{equation}

\paragraph{Self-consistency check.} For this approximation to hold, the biological premise $N_e s \gg 1$ must be satisfied when evaluating the implied selection coefficient:
\begin{equation}
    N_e s^*_{\text{strong}} = \frac{N_e V_A}{U_d} \gg 1
\end{equation}

For all wild vertebrate populations in our dataset, the observed $V_A$ and $N_e$ (derived from $\pi$) yield $N_e V_A / U_d \gg 1$. Therefore, the strong-selection limit is mathematically self-consistent for these data.

\paragraph{Connection to linkage decay ($Z$).} Substituting the equilibrium variance ($V_A \approx U_d s$) into the per-generation variance survival rate defined in the main text ($Z = 1 - V_m/V_A$) yields:

\begin{equation}
    Z = 1 - \frac{U_d s^2}{U_d s} = 1 - s \qquad (N_e s \gg 1)
\end{equation}
This establishes that under strong selection, the decay of statistical associations driving selective interference across the coalescent is governed directly by the selection coefficient (see also Buffalo and Kern 2024).

\paragraph{Biological Failure Mode.} While mathematically consistent, this limit is biologically irreconcilable with known sequence evolution. Evaluating Equation (\ref{eq:sHoC}) using pedigree-derived $V_A$ (typically 0.10–0.30) and plausible, yet conservative, vertebrate mutation rates ($U_d \approx 1$) necessitates an average selection coefficient of $s^*_{\text{strong}} \approx 0.10–0.30$. This implied value exceeds the mean heterozygous selection coefficient established from empirical vertebrate DFEs ($\bar{s} \approx 10^{-3}–10^{-2}$; Eyre-Walker and Keightley 2007; Huber et al. 2017) by one to two orders of magnitude. For pedigree $V_A$ to represent stable genetic architecture under this limit, the entire variance must be driven by semi-lethal mutations, contradicting genomic data.

\subsection{The Weak-Selection Boundary ($N_e s \ll 1$)}
Conversely, if selection is extremely weak relative to drift, alleles drift nearly neutrally and selective depletion ($sV_A$) becomes negligible compared to stochastic loss ($V_A/(2N_e)$). Equation (\ref{eq:VA_MSD}) reduces to:

\begin{equation}
    V_A \approx 2N_e U_d s^2 \qquad (N_e s \ll 1)
\end{equation}

Solving for the implied near-neutral selection coefficient:

\begin{equation}
    s^*_{\text{weak}} = \sqrt{\frac{V_A}{2N_e U_d}}
\label{eq:sneut}
\end{equation}

\paragraph{Self-consistency failure.} We test whether the assumption $N_e s \ll 1$ holds when parameterized by empirical data:

\begin{equation}
    N_e s^*_{\text{weak}} = \sqrt{\frac{N_e V_A}{2U_d}} \ll 1 \quad \Rightarrow \quad N_e \ll \frac{2U_d}{V_A}
\end{equation}

Using standard values ($V_A = 0.20$ and $U_d = 1$), this inequality demands an effective population size $N_e \ll 10$. Because all monitored populations in our dataset maintain sequence-derived $N_e$ orders of magnitude larger than this, the weak-selection limit collapses under these parameters. Attempting to fit the data to this regime yields $N_e s^*_{\text{weak}} \gg 1$, violating the foundational premise of near-neutral drift.

\subsection{Equivalence with Single-Locus Diffusion Theory}
To verify that the discrete-generation balance model above is robust, we map it to classical continuous-time diffusion theory. For a single locus, the stationary allele frequency reflects a flux balance between mutational influx, $\mu$, selective removal, $s \cdot \mathbb{E}[q]$, and stochastic loss due to genetic drift, $\mathbb{E}[q]/(2N_e)$. Setting this net flux to zero yields the expected allele frequency:

\begin{equation}
    \mu = \mathbb{E}[q]\left(s + \frac{1}{2N_e}\right) \quad \Rightarrow \quad \mathbb{E}[q] = \frac{\mu}{s + \frac{1}{2N_e}}
\end{equation}
(Crow and Kimura 1970). The per-locus additive variance is $2s^2 \mathbb{E}[q]$. Expanding this to the genome-wide variance by summing across $L = U_d/(2\mu)$ independent haploid functional targets:

\begin{equation}
    V_A = \left( \frac{U_d}{2\mu} \right) \left( 2s^2 \right) \left( \frac{\mu}{s + \frac{1}{2N_e}} \right) = \frac{U_d s^2}{s + \frac{1}{2N_e}} = \frac{2N_e U_d s^2}{2N_e s + 1}
\end{equation}

This recovers Equation (\ref{eq:VA_MSD}), confirming that the discrete-time derivation aligns perfectly with classical diffusion approximations under the assumption of rare, additive alleles.

\subsection{Generalization to Empirical Fitness Architectures (DFE)}
The preceding derivations assume a uniform selection coefficient, $s$, across all functional sites. However, empirical vertebrate architectures are governed by complex distributions of fitness effects (DFE). To evaluate the boundary conditions under biological architecture, we integrate across the DFE following the framework of Rettelbach et al. (2025).

\paragraph{Reconciling Pedigree $V_A$ with Empirical DFEs.}
Under House of Cards mutation-selection balance, the per-locus variance contribution of a deleterious mutation (with heterozygous disadvantage $t \equiv s$; we use the symbol $t$ to be consistent with the Rettelbach et al. derivations) at equilibrium frequency ($q \approx \mu/t$) is $v_A \approx 2\mu t$. Integrating this variance contribution across the probability density function of fitness effects, $\phi(t)$, yields the total genome-wide variance:

\begin{equation}
    V_A = U_d \int_0^\infty \phi(t), t, dt = U_d\bar{t}
\label{eq:VA_DFE}
\end{equation}

Solving for the implied mean selection coefficient required to support the observed $V_A$:

\begin{equation}
    \bar{t}^* = \frac{V_A}{U_d}
\label{eq:tbar_implied}
\end{equation}

This confirms that the single-locus House of Cards result ($V_A \approx U_d s$) holds for the arithmetic mean of the DFE. For the empirical $V_A \in [V_{A,\text{min}}, V_{A,\text{max}}]$ and $U_d \approx \text{[XX]}$ in our dataset, this relationship requires an average selection coefficient $\bar{t}^* \in [\text{XX}, \text{XX}]$. Because this vastly exceeds the empirical vertebrate DFE mean ($\bar{t} \approx 10^{-3}–10^{-2}$), this DFE-parameterized derivation provides an independent validation that pedigree-derived variance magnitudes cannot be maintained at mutation-selection-drift equilibrium.

\paragraph{Forward Prediction of Stable Variance.}
Alternatively, we can use established empirical DFE parameters to predict the maximum sustainable variance. Parameterizing the four-class uniform DFE of Rettelbach et al. (2025) and integrating only across classes governed by deterministic selection ($2N_e t > 1$):
\begin{equation}
    V_A^{\text{DFE}} = U_d \sum_{i=1}^{3} f_i \cdot \frac{t_{i+1} + t_i}{2}
\label{eq:VA_forward}
\end{equation}
where $f_i$ denotes the frequency of weakly, moderately, and strongly deleterious classes. Populating this with vertebrate-calibrated parameters yields $V_A^{\text{DFE}} \approx \text{[XX]}$. Because this theoretically sustainable variance is substantially lower than the contemporary $V_A$ estimated from pedigrees, the structural mismatch remains regardless of the assumed fitness architecture.

\paragraph{Strong-Selection Correction for Unlinked Interference.}
When variance is disproportionately driven by strongly deleterious mutations, or the "tail" of the DFE, the Robertson formulation slightly overestimates the severity of selective interference on unlinked neutral sites. Because strongly deleterious alleles are purged rapidly ($q \approx \mu/t \to 0$), they generate minimal multi-generational linkage disequilibrium despite their large instantaneous variance contribution.

The corrected unlinked interference factor derived from diffusion theory is:
\begin{equation}
    B_\text{unlinked} \approx \exp\left(-4U_d\int_0^\infty \phi(t)\frac{t}{(1+t)^2},dt\right) = \exp(-4U_d\bar{t}_{\text{eff}})
\label{eq:B_corrected}
\end{equation}

The effective mean driving this suppression is:

\begin{equation}
    \bar{t}_{\text{eff}} = \int_0^\infty \phi(t)\frac{t}{(1+t)^2}\,dt < \bar{t}
\end{equation}

Because $\bar{t}_{\text{eff}} < \bar{t}$, the true unlinked suppression ($B_\text{unlinked}^\text{corrected}$) is weaker than the baseline Robertson expectation ($B_\text{unlinked} = \exp(-4V_A)$). At the highly elevated $\bar{t}^*$ values implied by pedigree data, the fractional correction is approximately $1/(1+\bar{t}^*)^2 \in [\text{XX}, \text{XX}]$, meaning the Robertson formula overstates unlinked suppression by XX\%–XX\%. Therefore, the interference predictions generated in the main text are strictly conservative bounds: the true coalescent suppression at the implied high selection coefficients is weaker than predicted, and implementing this DFE-based correction fails to rescue the predicted $N_e$ back to empirically observed sequence diversity levels.

\end{document}